\chapter*{Introduction}
\addcontentsline{toc}{chapter}{Introduction}
\label{ch:introduction}

Objects of Research 

Subject of Research 

Aim OF The Study 

The Purpose of the Dissertation 

Tasks of Dissertation 

The Main Provisions to be Defended:
Degree of Reliability 

Approbation of the work 

Personal Contribution of the Author.
CHAPTER ONE \protect\hyperlink{_Toc215220436{18}}
\section{Conceptual Clarifications \protect\hyperlink{_Toc215220437{18}}}
\label{section:1-1}
\subsection{Concept of Urbanization}
\label{subsection:1-1-1}
\subsection{Urbanization in the Tropical Context}
\label{subsection:1-1-2}
\subsection{Land Use and Land Cover (LULC)}
\label{subsection:1-1-3}
\subsection{Urbanization Gradient}
\label{subsection:1-1-4}
\subsection{Plant Diversity}
\label{subsection:1-1-5}
\subsection{Ecosystem Functioning}
\label{subsection:1-1-6}
\subsection{Vegetation Indices (NDVI, EVI)}
\label{subsection:1-1-7}
\subsection{Soil Ecosystem Indicators}
\label{subsection:1-1-8}
\subsection{Urban Ecology}
\label{subsection:1-1-9}
\section{Review of Related Concepts and Approaches}
\label{section:1-2}
\subsection{Land Use / Land Cover Change Detection}
\label{subsection:1-2-1}
\subsection{Plant Diversity Assessment in Urban Environments}
\label{subsection:1-2-2}
\subsubsection{Vegetation Structure in Urban Landscapes}
\label{subsubsection:1-2-2-1}
\subsubsection{Biodiversity Indices}
\label{subsubsection:1-2-2-2}
\subsubsection{Statistical Analysis of Plant Diversity}
\label{subsubsection:1-2-2-3}
\subsection{Ecosystem Functioning Across Urban Gradients}
\label{subsection:1-2-3}
\subsection{Vegetation Indices in Urban Ecology Research}
\label{subsection:1-2-4}
\subsection{Soil Quality and Urbanization}
\label{subsection:1-2-5}
\subsection{Urbanization Metrics and Ecological Impacts}
\label{subsection:1-2-6}
\subsection{Statistical and Modelling Approaches in Urban Ecology}
\label{subsection:1-2-7}
\section{Theoretical Framework}
\label{section:1-3}
\subsection{Landscape Ecology Theory}
\label{subsection:1-3-1}
\subsection{Species Area Relationship (SAR) Theory}
\label{subsection:1-3-2}
\subsection{Intermediate Disturbance Hypothesis (IDH)}
\label{subsection:1-3-3}
\subsection{Urban Scaling Theory}
\label{subsection:1-3-4}
\subsection{Ecosystem Function Theory}
\label{subsection:1-3-5}
\section{Empirical Review}
\label{section:1-4}
\subsection{Global Studies}
\label{subsection:1-4-1}
\subsection{African Studies}
\label{subsection:1-4-2}
\subsection{Nigerian Studies}
\label{subsection:1-4-3}
\section{Summary of Research Gaps}
\label{section:1-5}
\section{Conceptual Framework}
\label{section:1-6}

CHAPTER TWO 
\section{RESEARCH METHODOLOGY}
\label{section:2-0}
\section{Research Design}
\label{section:2-1}
\section{Study Area}
\label{section:2-2}
\section{Data Requirements and Sources}
\label{section:2-3}
\subsection{Remote Sensing Data}
\label{subsection:2-3-1}
\subsection{Field Vegetation Data}
\label{subsection:2-3-2}
\subsection{Soil Data}
\label{subsection:2-3-3}
\section{Land-Use and Land-Cover (LULC) Analysis}
\label{section:2-4}
\subsection{Image Pre-Processing}
\label{subsection:2-4-1}
\subsection{LULC Classification Using Random Forest}
\label{subsection:2-4-2}
\subsection{Accuracy Assessment}
\label{subsection:2-4-3}
\subsection{LULC Change Detection}
\label{subsection:2-4-4}
\section{Plant Diversity Assessment}
\label{section:2-5}
\subsection{Sampling Design}
\label{subsection:2-5-1}
\subsection{Biodiversity Indices Computation}
\label{subsection:2-5-2}
\subsection{Statistical Analysis of Diversity}
\label{subsection:2-5-3}
\subsubsection{One-Way ANOVA}
\label{subsubsection:2-5-3-1}
\subsubsection{Post-hoc Tests}
\label{subsubsection:2-5-3-2}
\section{Ecosystem Functioning and Urbanization Impact Modelling}
\label{section:2-6}
\subsection{Urbanization Metrics Extraction}
\label{subsection:2-6-1}
\subsection{Ecosystem Function Indicators}
\label{subsection:2-6-2}
\subsection{Zone Comparison Statistics}
\label{subsection:2-6-3}
\subsection{Regression Modelling}
\label{subsection:2-6-4}
\section{Policy and Management Recommendations}
\label{section:2-7}

CHAPTER THREE 

RESULTS AND DISCUSSION 
\section{Land Use and Land Cover (LULC) Classification Results}
\label{section:3-1}
\subsection{Land Use and Land Cover Dynamics}
\label{subsection:3-1-1}
\subsection{Spatial Interpretation of LULC Maps}
\label{subsection:3-1-2}
\subsection{Ground Truthing for LULC Classification}
\label{subsection:3-1-3}
\section{Plant Composition and Biodiversity Patterns}
\label{section:3-2}
\subsection{Plant Diversity in Makurdi local Government area}
\label{subsection:3-2-1}
\subsubsection{Rural Areas of Makurdi Local Government}
\label{subsubsection:3-2-1-1}
\subsubsection{Peri-Urban Areas of Makurdi Local Government}
\label{subsubsection:3-2-1-2}
\subsubsection{Urban Areas of Makurdi Local Government}
\label{subsubsection:3-2-1-3}
\subsubsection{Comparison of Species Diversity between Rural, Peri-Urban, and Urban Areas in Makurdi Local Government}
\label{subsubsection:3-2-1-4}
\subsection{Plant Diversity in Otukpo Local Government}
\label{subsection:3-2-2}
\subsubsection{Peri-Urban Area of Otukpo Local Government}
\label{subsubsection:3-2-2-2}
\subsubsection{Rural Areas of Otukpo Local Government}
\label{subsubsection:3-2-2-1}
\subsubsection{Urban Areas of Otukpo Local Government}
\label{subsubsection:3-2-2-3}
\subsubsection{Comparison of Species Diversity between Rural, Peri-Urban, and Urban Areas in Otukpo Local Government}
\label{subsubsection:3-2-2-4}
\section{Soil Properties Across the Urban Gradient}
\label{section:3-3}
\subsection{Soil Variable Summary}
\label{subsection:3-3-1}
\subsection{Statistical Differences Across Zones}
\label{subsection:3-3-2}
\section{Vegetation Index (NDVI/EVI) Results}
\label{section:3-4}
\subsection{NDVI and EVI Summary by Zone}
\label{subsection:3-4-1}
\subsection{Zone Comparison Statistics}
\label{subsection:3-4-2}
\section{Relationship Between Urbanization Metrics and Ecosystem Functioning}
\label{section:3-5}

3. 5.1 Summary of Urbanization Indicators 
\subsection{Regression Analysis of Urbanization Metrics as Predictors of NDVI}
\label{subsection:3-5-2}
\subsection{Model Fit Validation}
\label{subsection:3-5-3}
\subsection{Interpretation of Regression Results}
\label{subsection:3-5-4}
\section{Integrated Assessment of Urbanization Impacts}
\label{section:3-6}
\subsection{Synthesis of Findings Across Indicators}
\label{subsection:3-6-1}
\subsection{Ecological Implications}
\label{subsection:3-6-2}
\section{Summary of Findings}
\label{section:3-7}
\section{Discussion of Results}
\label{section:3-8}
\subsection{Urban Expansion and Landscape Transformation}
\label{subsection:3-8-1}
\subsection{Plant Diversity Decline Along Urban Gradients}
\label{subsection:3-8-2}
\subsection{Soil Properties and Urban-Induced Alterations}
\label{subsection:3-8-3}
\subsection{Vegetation Productivity Decline (NDVI/EVI)}
\label{subsection:3-8-4}
\subsection{Influence of Urbanization Metrics on Ecosystem Functioning}
\label{subsection:3-8-5}
\subsection{Synthesis with Tropical Urban Ecology Theory}
\label{subsection:3-8-6}

CONCLUSION 

Conclusion 

RECOMMENDATION 

Policy Implications 

Contribution to Knowledge 

List of publications on the topic of the thesis 

REFERNCES \protect\hyperlink{_Toc215220540{110}}

LIST OF ABBREVIATIONS 

APPENDICES 

{}\textbf{INTRODUCTION}

\textbf{Relevance of the Study.} Urbanization can be seen as a major cause of ecological change in tropical regions, and still its impacts on biodiversity and ecosystem functioning remain understudied in many rapidly growing African cities. To understand how expanding built-up areas alter plant communities, soil processes, and vegetation productivity is important for detecting long-term environmental consequences (Obateru et al. 2025). In recent studies (Obateru et al. 2025; Chen et al. 2023; Mhlanga et al. 2025), it is evident that tropical cities in particularly faces intense pressures as a result of high population growth, weak land-use regulation, and high dependence on natural resources, making them global hotspots for biodiversity decline and habitat fragmentation. To address this gap, this study provides spatial, ecological, and statistical evidence of how urbanization affects plant diversity and ecosystem functioning in Benue State, Nigeria\textbf{.}

This study is particularly relevant because it brings together remote sensing, field-based biodiversity surveys, and soil ecosystem indicators an integrated approach that is still rare in West African ecological research. Most previous studies have focused either on satellite-derived vegetation analysis or on field species surveys, often treating them separately and rarely examining urban, peri-urban, and rural areas together (Onandia et al., 2019; Li et al., 2016). To address this gap, this study adopts a more holistic view of ecological health, capturing how land-use change interacts with vegetation structure and soil nutrient dynamics (Finizio et al., 2024).
This study is also important for urban planning and environmental policy. Many Nigerian cities, including Makurdi and Otukpo, lack up-to-date ecological information to guide land-use decisions, green-space protection, and climate adaptation efforts. In their respective studies, (Liu et al., 2025; Akiner \& Ghasri, 2025) shows that integrating biodiversity indicators and vegetation metrics such as NDVI into planning processes can greatly enhance urban sustainability. The outcome of this research will provide the policymakers practical, evidence-based guidance for developing green infrastructure, conserving remaining natural habitats, and supporting ecosystem services that are essential for human well-being.
Finally, this study adds to global discussions on how urbanization influences ecological processes, especially in data-limited regions. By developing a statistical model that links urbanization intensity measured through built-up density, night-time lights, and population density with ecosystem indicators such as NDVI, soil organic carbon, and nitrogen, the research improves understanding of urban ecological gradients in tropical settings. Again, for researchers studying similar urban transitions across Sub-Saharan Africa and other rapidly urbanizing tropical regions the study finding will provide them with useful comparisons (Stenchly et al., 2019; Ayeni et al., 2023). In doing so, the study not only addresses regional data gaps but also strengthens broader ecological research and practice.
\textbf{The Degree of Development of the Problem.} Although urban ecological issues are receiving increasing global attention, past literatures on how urbanization affects plant diversity and ecosystem functioning in tropical African cities remains relatively scarce. Existing studies have largely concentrated on temperate regions, where long-term ecological monitoring systems and high-resolution datasets are more readily available (Rija et al. 2014). In contrast, cities in tropical Africa particularly in West Africa are experiencing rapid and often poorly regulated land-use changes, yet they are far less represented in the scientific literature. The overall knowledge gaps are continues as regarding the impacts of accelerated urban growth on local plant communities, soil conditions, and vegetation productivity along rural--urban gradients (Radoine et al. 2024; Atchade et al. 2023).
The literatures that explicitly connects these parts; urbanization, biodiversity, and ecosystem functioning in Nigeria are rare. Most existing studies tend to focus on seperate components like LULC, plant diversity, or soil degradation without examining their combined or interactive effects within an integrated framework (Ayeni et al. 2023). In the other hand, some of the literatures that apply quantitative urbanization indicators, to include night-time light intensity, built-up density, and population concentration, remain particularly limited, despite growing evidence that these metrics are strong predictors of ecological stress at the global scale (O'Connor et al. 2022). This fragmented approach has constrained a holistic understanding of how urban expansion is reshaping ecological processes in rapidly growing cities such as Makurdi and Otukpo.
The problem is also underdeveloped in terms of methodological integration. While remote sensing technologies have improved substantially, their application to tropical urban ecology is still limited by insufficient ground-truthing and the absence of fine-scale biodiversity data (Hisano, et al. 2017). Likewise, in Nigeria literatures centred on plant diversity are often short-term and restricted to small sampling areas, lacking the extensive spatial stratification required to capture the urban gradients (Gao and Carmel 2020). Also, Soil indicators such as organic carbon and nitrogen that are vital in the evaluation of ecosystem functioning, are rarely incorporated into ecological impact assessments despite mounting evidence of their sensitivity to urban disturbance (Udoinyang 2024).
The conceptual and empirical understanding of this issue remains incomplete. Because only a limited number of studies have applied statistical modeling approaches to quantify how multiple dimensions of urbanization jointly influence vegetation health, soil conditions, and plant community structure in tropical settings. In the most recent innovation in the filed of ecology that emphasize the pressing need for integrating, data-driven frameworks that combine remote sensing, species inventories, and ecological indicators to better capture the complex impacts of urbanization (Hines and Pereira 2021; Ferreira et al. 2019). By responding to these gaps, the present study makes a meaningful contribution toward advancing research on this largely understudied problem in tropical cities---particularly within Benue State and Nigeria more broadly.
Understanding the impacts of urbanisation and its challenges on natural ecosystems is important. The role it plays in shaping the ecological and economic sphere in Benue state and nigera at large has led to different research by several scientific studies ( Shabu.T, Garba U., Job l.o,Mala.M, Mabogunje A.L, Onokerhoraye A.G, Eghaosa E.O, Ismaila.T, Guessan.K.N, Jelili.M.O,Ajibade A., Alabi A.T., Stren R.E, Faniran.A. Onanuga, M. Y.,Nnaji, C. C, Tyubee, B. T., Onilude, O,Ogunbode, T.
{}\textbf{Objects of Research}


We study and assess the effects of urbanisation on plant diversity and ecosystem services in Makurdi and Otukpo, two rapidly urbanising cities in Benue State, Nigeria. This study examine LULC over three decade ranging from 1990-2024 to understand the impact of urban expansion on local plants.
{}\textbf{Subject of Research}

Multi-temporal satellite imagery and ancillary geospatial datasets were employed which are necessary for detecting land-use/land-cover transformations, quantifying vegetation condition, and assessing urbanization intensity across Makurdi and Otukpo. Landsat TM (1990), ETM+ (2010), and OLI (2024) scenes were obtained for change detection, NDVI, and EVI derivation due to their 30 m spatial resolution and long-term continuity, which makes them suitable for landscape-scale ecological studies Downloading Landsat imagery that covers both Makurdi and~~~ Otukpo, at interval~~ of 10 years between 1990, 2010, and 2024 years for possibly three period. VIIRS Day/Night Band night-time lights imagery was used to represent anthropogenic intensity and built-up expansion through illumination patterns, while WorldPop population density rasters provided high-resolution demographic information to quantify population-driven urbanization pressures these were extracted at the ward level.
{}\textbf{Aim OF The Study}

The major aim and purpose of this study is to investigate the impacts of urbanization on plant diversity and ecosystem functioning in Benue State, with a focus on the changes in LULC patterns, and the resulting effects on ecosystem services and plants biodiversity. The selected time periods for analysis are 1990, 2010, and 2024, reflecting significant phases of urbanization and environmental change in the study area.
{}\textbf{The Purpose of the Dissertation}

To assess how urbanization influences plant diversity and ecosystem functioning across urbanization gradients in Makurdi and Otukpo LGAs Benue state using an integrated approach combining remote sensing, field-based biodiversity assessment, soil indicators, and statistical modelling.
To analyze the changes in LULC in Makurdi and Otukpo for 1990, 2010, and 2024, and to quantify the trends and extent of urban expansion and vegetation loss across rural, peri-urban, and urban zones.
To assess plant species richness, diversity, and evenness across urban gradients using field-based vegetation inventory and compute indices such as Shannon index, Simpson index, and Pielou's evenness.
To evaluate how increasing urbanization influences the ecosystem functioning across rural-urban gradients, and statistically determine whether urbanization significantly affects plant diversity and soil ecosystem quality.
To develop and propose an evidence-based policy recommendations for sustainable urban planning and biodiversity conservation in tropical cities, based on the study's findings on urbanization impacts on plants and ecosystem function.
Scientific Novelty.
\begin{itemize}
\item
  The systematic sampling of 33 vegetation plots across all wards in Makurdi and Otukpo further provides one of the few fine-scale biodiversity datasets for the region, strengthening empirical understanding of species-environment interactions in tropical cities.
\item
  The study contributes to global scientific debates on how urbanization shapes ecological processes, particularly in regions where data scarcity limits robust analysis.
\item
  This study integrates field based plant diversity assessments, soil ecosystem indicators, and remote sensing metrics to evaluate urbanization impacts in a tropical city, a combination rarely applied in Nigerian or wider tropical urban ecological research.
\item
  Through a quantitative regression framework, the research introduces a more robust, multi-dimensional assessment of anthropogenic pressure an approach underutilized in ecological studies in the research area, by Connecting NDVI to built-up density, night-time lights, and population density
\item
  The study methodological and empirical contributions advances current knowledge on how rapid tropical urbanization shapes plant diversity and ecosystem functioning, addressing well-recognized research gaps in urban ecology literature.
\end{itemize}

{}\textbf{Tasks of Dissertation}

Remote sensing served as the foundation for quantifying land-use transitions for a selected period from 1990, 2010, 2024 and vegetation dynamics, using satellite-derived indices and classification products processed within the Google Earth Engine (GEE) environment. The Field based plants sampling was carried out across 33 stratified plots, representing urbanization gradientsin the 6 wards of Makurdi and 5 wards of Otukpo (three plots per ward. For each plot, GPS coordinates, elevation, land-use context, and additional plot metadata were documented to ensure accurate geolocation and integration with remotely sensed variables. These data supported the computation of biodiversity indices (Shannon, Simpson, Evenness) and enabled comparison of species composition along the rural--peri-urban--urban gradient. Complementary to the plant survey, geo-cordinates of the plotss were recorded.
Statistical analysis, including ANOVA tests for comparing vegetation and soil parameters across the urbanization gradient, with assumption checks for data distribution and homogeneity. The Regression modelling was implemented to assess the relationship between urbanization metrics and vegetation condition, using the specified model. Vegetation indices results showed the clearest ecological signal. NDVI declined from 0.482 ± 0.023 (rural) to 0.277 ± 0.076 (urban), while EVI dropped from 0.312 ± 0.025 to 0.173 ± 0.051, demonstrating stark reductions in primary productivity and vegetation vigor. ANOVA results confirmed highly significant differences exist across the gradients zones (NDVI: F = 31.04, p \textless{} 0.001; EVI: F = 32.52, p \textless{} 0.001), revealing that vegetation functioning diminishes sharply with increasing urban intensity. Ecosystem functioning indicators showed a marked decline with increasing urban intensity. The NDVI values have decreased from 0.48 in the rural areas to 0.28 in the urban areas, while the EVI values also have reduced from 0.31 to 0.17, which indicate the reduction of vegetation vigor and the primary productivity. On the other hand, the soil analysis\textquotesingle s findings showed that a significant increases in the soil pH in urban zones is (p = 0.037), which suggested there is an altered soil chemistry linked to anthropogenic disturbance. Although SOC and nitrogen appeared higher in urban zones, these increases are likely driven by waste inputs rather than natural fertility processes. On the other hand, the regression modelling further showed that the population density is the strongest suppressor of vegetation greenness having this value (β = --0.0072, p \textless{} 0.001), highlighting demographic pressure as a key driver of ecological degradation.
To accurately know if vegetation indices differ significantly across the urban zones, one-way ANOVA tests were conducted for NDVI and EVI. The findings indicate highly significant differences where (NDVI: F = 31.04, p \textless{} 0.001, η² = 0.75; EVI: F = 32.52, p \textless{} 0.001, η² = 0.76), proving the impact urbanization has on plant diversity.
The overall model performance (R² = 0.862) shows NDVI variability is accounted for by these three indicators, underscoring the central role of urbanization dynamics in shaping ecological conditions. The strong negative effect of population density combined with the modest positive effect of built-up surfaces suggests that ecological degradation is predominantly linked to demographic expansion rather than structural development alone.
Biodiversity Decline and Ecological Homogenization.The decline in species richness from 23.50 species (rural) to 15.83 species (urban), indicates that urbanization fosters homogenization of plant communities. Disturbance-tolerant species such as Cleomeviscosa, Sida acuta, and Heterotis rotundifolia become more dominant, while native forest species decline. This shift reduces functional diversity, disrupts trophic and symbiotic relationships, and diminishes ecological resilience.
Habitat Fragmentation and Landscape Conversion: The loss of 124.93 km² of forest and widespread cropland decline have fragmented formerly continuous habitats into isolated patches. Fragmentation reduces dispersal opportunities, increases edge effects, and disrupts plant community dynamics.
\textbf{Theoretical and Practical Significance of the Work.} From a theoretical perspective, the study deepens understanding of the effect of urbanization on ecological patterns in tropical cities (urban ecology, landscape ecology, and socio-environmental systems). It advances existing theory by empirically connecting key measures of urbanization such as built-up density, night-time light intensity, and population density with indicators of ecosystem functioning, including vegetation greenness (NDVI), soil nutrients, and plant diversity. The study demonstrated how urban growth and the ecosystem can be modelled (Reynolds 2002; Zaninotto et al. 2023; Ulrich and Sargent 2025). The findings provides new insights for resilience, species distribution patterns, and the pathways for ecosystem processes (Xie et al. 2020; Isendahl et al. 2020).
From a practical view, the results provide actionable insights for urban planners, environmental managers, and policymakers in Benue State and other tropical regions experiencing rapid urban growth. The high-resolution spatial and field datasets produced including vegetation surveys, soil indicators, and land-use/land-cover change maps can be directly applied to environmental monitoring, impact assessments, and local climate adaptation planning. Lastly, the findings identified priority areas for conservation and ecological restoration also proven in past studies (Korpilo et al. 2023; Hansen et al. 2023) this enables targeted interventions that reduce biodiversity loss and the enhancement of ecosystem services in mostly growing and expanding cities (Falasca and Marucci 2024; Jensen et al. 2025).
\textbf{Methodology and Methods of Research.} This study was carried out between 2022 and 2025 across the two study loactions. The research design follows an integrated methodological framework that was structured around the major study objectives, and by combining geospatial analysis, ecological field measurements, and soil assessments. Remote sensing served as the foundation for quantifying land-use transitions for a selected period from 1990, 2010, 2024 and vegetation dynamics, using satellite-derived indices and classification products processed within the Google Earth Engine (GEE) environment. Field based vegetation sampling complemented these spatial datasets by documenting species composition and diversity across gradients of urbanization. Soil properties essential for understanding ecosystem functioning such as Soil Organic Carbon, Nitrogen, Bulk Density, pH, and Cation Exchange Capacity were extracted using GEE and also the urbanization indicators were etracted from the same GEE. The methodology used provides a comprehensive, multi-dimensional approach for achieving the purpose of the study.
{}\textbf{The Main Provisions to be Defended:}

The study examines how urban expansion, vegetation transformation, and anthropogenic pressure over a 34-year period (1990-2024) shape ecological patterns across rural, peri-urban, and urban zones. The principal scientific provisions for defense in this work are summarized as follows:
Regression Analysis of Urbanization Metrics as Predictors of NDVI.
The model performed strongly, indicating that a considerable proportion of spatial variation in NDVI is attributable to these three urbanization indicators. The overall model explained 86.2\% of the variation in NDVI (R² = 0.862), with an adjusted R² of 0.842, confirming high explanatory power even after adjusting for the number of predictors.The model's F-statistic was significant where (F = 41.73, p \textless{} 0.001) indicating that the predictor set collectively contributes meaningfully to NDVI variation. Information criteria (AIC = --80.28; BIC = --75.57) also support a well-fitted model.
Synthesis of Findings Across Indicators

Vegetation indices showed the clearest ecological signal. The value (NDVI) declined from 0.482 ± 0.023 to 0.277 ± 0.076, while EVI dropped from 0.312 ± 0.025 to 0.173 ± 0.051, demonstrating stark reductions in primary productivity and vegetation vigor. ANOVA results confirmed highly significant differences across zones (F = 31.04, p \textless{} 0.001 in the case of NDVI; and F = 32.52, p \textless{} 0.001 for the case of EVI), revealing that vegetation functioning diminishes sharply with increasing urban intensity

Habitat Fragmentation and Landscape Conversion

The loss of 124.93 km² of forest land and also the drastical reduction in cropland size have fragmented the formerly continuous habitats into isolated patches. This fragmentation in turn reduces dispersal opportunities, increases edge effects, and disrupts plant community dynamics. And again, patch isolation also limits pollinator movement, seed flow, and successional processes essential for ecosystem recovery.
Biodiversity Decline and Ecological Homogenization

The decline in species richness from 23.50 species (rural) to 15.83 species (urban), along with reduced Shannon and Simpson indices, indicates that urbanization fosters homogenization of plant communities. Disturbance tolerant species such as Cleome viscosa, Sida acuta, and Heterotis rotundifolia become more dominant, while native forest species decline. The shift reduced functional diversity, disrupts trophic and symbiotic relationships, and diminishes ecological resilience.
Soil Fertility Modification and Biological Function Disruption

Significantly higher soil pH in urban areas (p = 0.037) indicates altered nutrient availability and potential micro-nutrient deficiencies that can suppress plant growth. Even though SOC and nitrogen appear elevated, the biological reality is reduced nutrient turnover, poorer microbial activity, and compacted soil structure conditions that limit root development and reduce soil-plant interactions fundamental to ecosystem processes.
Reduced Primary Productivity and Vegetation Health

The substantial decline in NDVI value from 0.48 to 0.28 and that of EVI from 0.31 to 0.17 tells the weakened photosynthetic activity, reduced biomass accumulation, and lower canopy cover. As a result of the reduction it compromises ecological functions such as carbon sequestration, microclimate regulation, soil stabilization, and hydrological buffering. Urban hotpots with the highest population densities show the lowest NDVI values, confirming that demographic stress is a major biological driver of ecosystem decline.
Statistical Models Confirm Urbanization as a Predictor of Ecological Degradation.
The regression modeling of these variables (NDVI = β₀ + β₁ Built-up + β₂ Night-time Lights + β₃ Population Density + ε) revealed that the three indicators of urbanization had significantly predicted plant productivity in the study area. Random Forest regression further ranked the predictors by importance, confirming built-up area and night-time lights as the leading determinants of ecosystem functioning in both LGAs.
Evidence for Sustainable Urban Planning and Conservation Interventions.
The combined findings from LULC change, biodiversity patterns, soil indicators, and statistical modeling support actionable recommendations for urban management in Benue State. Key provisions include:
\begin{itemize}
\item
  Strengthening land-use zoning and control of built-up expansion.
\item
  Restoration of degraded vegetation and promotion of urban forestry.
\item
  Protection of remaining natural vegetation patches in peri-urban and rural zones.
\item
  Integrating ecological datasets (NDVI, species diversity, soil metrics) into planning decisions.
\end{itemize}

These provisions align with global frameworks for nature-based solutions and provide scientific justification for policies that enhance biodiversity conservation and ecosystem resilience in tropical urban environments.
{}\textbf{Degree of Reliability}

All the classified maps were validated with ground truth data collected from field through accuracy assestment tool in ArcGIS. A total of 11 wards were randomly picked (6 in Markurdi and 5 in Otukpo), where each ward had 3 plot samples and each plot sample covered an area of 100 m². All were surveyed in each zone. The plant species were manually identified, and classified based on their life forms (trees, shrubs, and herbs) and growth forms (perennial or annual). Species richness and abundance were also recorded for each plot. This survey was carried out alongside the Department of Wildlife and Range Management, College of Forestry and Fisheries, Joseph Sarwuan Tarka University, Makurdi. All geospatial predictor variables used in this study including Enhanced Vegetation Index (EVI), Normalized Difference Vegetation Index (NDVI), soil organic carbon, nitrogen, bulk density, cation exchange capacity, pH, night-time lights, population density, and built-up intensity were extracted using Google Earth Engine (GEE)(Kumari et al. 2021; Hengl et al. 2017; Song et al. 2025; Wang et al. 2024).
The workflow integrated multiple remotely sensed and modelled datasetsand all layers were clipped to ward-level buffer boundaries, aggregated to mean values per ward, and exported directly from GEE into a structured table. The final datasets produced by this workflow is openly available in the Appendix, ensuring reproducibility as well as the national population data in appendix 10 was collected from the National Population Commission, Department of Population Census and Statistics, Headquarters, Abuja.
Compliance of the Dissertation with the Passport of the Scientific Specialty.
Dissertation work corresponds to the specialty passport 1.5.15. Ecology by item: Item 10. Anthropogenic impact on populations, communities, and ecosystems. Biological effects of environmental pollution with toxic substances (ecotoxicology). Development of biological methods and criteria for assessing the state of the environment, bioindication, biotesting, and biomonitoring. Development of environmentally sound standards for the impact of human economic activity on wildlife: item 12. Ecological principles of nature protection at the population-specific and ecosystem level.
{}\textbf{Approbation of the work}

The result of the dissertation research was the publication of 5 articles, including 3 in scientific journals indexed in the SCOPUS database and 2 in a journal included in the Higher Attestation Commission (HAC) list. The results of the dissertation work were presented at the following conferences: \textbf{International youth scientific conference ``EPI 2025''}: Environmental studies and protection issues, Institute of Environmental Engineering RUDN University April 2025, \textbf{XXVth International Multidisciplinary Scientific GeoConference Surveying, Geology and Mining, Ecology and Management (SGEM) A}lbena Resort \& Spa, BulgariaJune 28-July 7th 2025, proceeding, \textbf{Roles of Artificial intelligence in predictive Agroecological modelling and sustainable land management in response to climate viability in Moscow Region (Scopus Q2)}

Africa's pathway to the 17 SDGS: \textbf{An Integrated Approach November 21st, 2023: Abundant Potential Solar Energy in Nigeria: A Pathway to Renewable Energy, Russia African Conference, July 2023} The Application of Remote Sensing Techniques in Monitoring and Assessment of Inorganic Pollutants in Africa, Anthropogenic Activities: \textbf{A Threat to Natural Resources Sustainability Possible Solution (Council of Young scientists of the institute of ecology of RUDN University ``Lecture Seminars on Studying the Earth'' December, 3 \& 10, 2022}), Urbanization Challenges and Its Pressure On Natural Resources Sustainability In Nigeria \textbf{International youth scientific conference ``GREEN-2022'' Graduate Research in Ecology, Engineering and Nature; November 2022 Institute of Environmental Engineering RUDN University,} International youth scientific conference ``EPI 2024'' Environmental studies and protection issues. Institute of Environmental Engineering RUDN University "Common Natural Ecosystem Challenges in Nigeria": \textbf{Assessment of urbanization\textquotesingle s environmental impact in Makurdi Metropolis, Benue State, Nigeria. International Journal of Environmental Impacts}, Vol. 8, No. 2, pp. 393-400. https://doi.org/10.18280/ijei.080218 (Scopus Q3, A social approach to assessing the impact of urbanization on natural resources: A case study of Makurdi-Benue State, Nigeria) \textbf{International Research Journal №6 (156). URL: https://doi.org/10.60797/IRJ.2025.156.82 Issue: № 6 (156), 2025 (VAK K1(SCOPUS Q2, WOS Q1,} Spatiotemporal Analysis of Land Use and Land Cover Changes Driven by Urbanization in Makurdi, Nigeria (1990 to 2024): Implications for Cropland Loss, Environmental Risk, and Sustainable Urban Development (SCOPUS Q3 WOS Q3), Journal of Applied Environmental Research, Socio-ecological consequences of urban growth in dutsen kura gwari, Niger state, Nigeria, Soil fertility management practices among smallholder farmers in Ughelli North, Delta State, Nigeria. (2025). Tropical Agriculture, 102(3), 358--367.
{}\textbf{Personal Contribution of the Author.}

The candidate participated in setting the aim and specific objectives of the study, collected and analyzed the obtained material, processed and interpreted the data, and prepared publications in co-authorship.
Structure and Volume of the Dissertation. The thesis consists of an Introduction, 3 chapters, conclusions, Recommendations, List of Abbreviations, list of references, 11 appendices. The thesis is summarized and presented on 141 pages; the main text includes 15 figures and 21 tables.
Acknowledgement

My special thanks go to God for all grace and provision, to my scientific supervisor Ass. Prof. Pinaev Vladimir Evgenievich, who supervised this work; Prof. Savenkova Elena Viktorovna, Director of the Institute of Environmental Engineering at RUDN University; Prof. Kapralova Daria Olegovna, PhD coordinator at the Institute of Environmental Engineering at RUDN University; Prof. M.M. Redina; and all staff of the Institute of Environmental Engineering at RUDN University for their support and tutelage throughout the period of my academic pursuit at the institute. To my family, my wife and daughter, Mrs. Patience and Rehoboth Ebogocho Shaibu; my mother, Aladi Shaibu; siblings, etc.; Mr. John Nyiesue Ikyobo, principal technical officer 1, Department of Wildlife and Range Management; and Surv. Idoko F. Olechi, Department of Physical Planning: Geoinformatics Unit, Joseph Sarwuan Tarka University, Makurdi-Nigeria, for their invaluable assistance during collection of data for the research work on the topic of the dissertation.
This chapter reviews relevant literature that provides the scientific foundation for assessing how urbanization influences plant diversity and ecosystem functioning in tropical environments. It is organized around the major themes of the study, beginning with the concept, patterns, and drivers of urbanization, followed by documented effects of urban growth on plant species richness, evenness, and community structure along rural-urban gradients. The chapter also examines ecosystem functioning metrics such as vegetation indices (NDVI/EVI), soil organic carbon, and soil nitrogen and their ecological significance, alongside remote sensing and GIS techniques used in monitoring land-use change and vegetation dynamics. Furthermore, the review highlights statistical modelling approaches, including regression and machine-learning tools, that help quantify the relationships between urban intensity and ecological variables. The chapter concludes with a conceptual and theoretical framework that integrates these themes and guides the analytical direction of the present study.
\hypertarget{chapter-one}{%
